\subsection{Comparison of DELT-Hit and DELi}

Table~\ref{tab:delt-deli-comparison} summarizes the main workflow-level similarities and differences between DELT-Hit and DELi.

\newcommand{\cmark}{\checkmark}
\begin{longtable}{p{4.9cm}p{3.7cm}p{3.7cm}}
\caption{Comparison of DELT-Hit and DELi across major DEL analysis tasks.}
\label{tab:delt-deli-comparison} \\
\toprule
\textbf{Feature} & \textbf{DELT-Hit} & \textbf{DELi} \\
\midrule
\endfirsthead

\toprule
\textbf{Feature} & \textbf{DELT-Hit} & \textbf{DELi} \\
\midrule
\endhead

Input files & Excel workbook compiled to \texttt{config.yaml} & Library JSON, building-block CSV files, and decode-settings YAML \\
User-friendliness & Single workbook keeps library design, barcode definitions, and analyses in one file & More modular and explicit, but usually requires editing several structured input files \\
Raw read decoding & \cmark & \cmark \\
Error-tolerant barcode matching & \cmark & \cmark \\
Library enumeration & \cmark & \cmark \\
Enrichment / analysis methods & Counts-based enrichment and \texttt{edgeR} & MLE, NSC/\texttt{SD\_min}, normalized z-score, log-space z-score, PolyO, overlap analyses \\
Analysis reports & Organized analysis folders and static outputs & Automated HTML report with figures and CSV exports \\
Dual-Display & \cmark &  \\
Counts-driven filtered enumeration & \cmark &  \\
Interactive count-table dashboard & \cmark &  \\
Single-compound structure lookup &  & \cmark \\
Molecular property and representation generation & \cmark &  \\
\bottomrule
\end{longtable}

The synthetic demultiplex benchmark documented in the shared runbook uses generated FASTQ datasets to compare DELT-Hit and DELi on the canonical 2-, 3-, and 4-cycle libraries with 10 building blocks per cycle across 1M, 10M, 100M, and 1000M reads, and additionally includes an extended 2-cycle runtime sweep with 10, 100, and 1000 building blocks per cycle. Benchmarking was completed on a Slurm cluster with 12 CPUs and 128\,GB RAM allocated per job. The reported outputs focus on end-to-end runtime and memory consumption for the shared benchmark setup; the DELi 1000M-read runs exited with an out-of-memory state. Figure~\ref{fig:delt-deli-runtime-cycles} summarizes cycle-wise runtime scaling, Figure~\ref{fig:delt-deli-runtime-2cycle-bbpc} shows the extended 2-cycle runtime sweep, and Figure~\ref{fig:delt-deli-memory-cycles} summarizes peak RSS scaling.

\begin{figure}[htbp]
\centering
\includegraphics[width=0.32\textwidth]{\benchmarkroot/benchmarks/demultiplex/plots/synthetic_2cycle_10bbpc_runtime.png}\hfill
\includegraphics[width=0.32\textwidth]{\benchmarkroot/benchmarks/demultiplex/plots/synthetic_3cycle_10bbpc_runtime.png}\hfill
\includegraphics[width=0.32\textwidth]{\benchmarkroot/benchmarks/demultiplex/plots/synthetic_4cycle_10bbpc_runtime.png}
\caption{Cycle-wise runtime scaling for the canonical synthetic demultiplex benchmark at 10 building blocks per cycle across 1M to 1000M reads. The DELi 1000M-read runs exited with an out-of-memory state.}
\label{fig:delt-deli-runtime-cycles}
\end{figure}

\begin{figure}[htbp]
\centering
\includegraphics[width=0.75\textwidth]{\benchmarkroot/benchmarks/demultiplex/plots/synthetic_2cycle_bbpc_comparison_runtime.png}
\caption{Runtime scaling for the extended 2-cycle synthetic benchmark across 10, 100, and 1000 building blocks per cycle.}
\label{fig:delt-deli-runtime-2cycle-bbpc}
\end{figure}

\begin{figure}[htbp]
\centering
\includegraphics[width=0.32\textwidth]{\benchmarkroot/benchmarks/demultiplex/plots/synthetic_2cycle_10bbpc_peak_rss.png}\hfill
\includegraphics[width=0.32\textwidth]{\benchmarkroot/benchmarks/demultiplex/plots/synthetic_3cycle_10bbpc_peak_rss.png}\hfill
\includegraphics[width=0.32\textwidth]{\benchmarkroot/benchmarks/demultiplex/plots/synthetic_4cycle_10bbpc_peak_rss.png}
\caption{Cycle-wise peak RSS scaling for the canonical synthetic demultiplex benchmark at 10 building blocks per cycle across 1M to 1000M reads. The DELi 1000M-read runs exited with an out-of-memory state.}
\label{fig:delt-deli-memory-cycles}
\end{figure}

Within this benchmark setup, DELT-Hit demonstrates more favorable run-time and memory complexity than DELi.
